%!TEX root = ../thesis.tex
%*******************************************************************************
%*********************************** First Chapter *****************************
%*******************************************************************************

\chapter{Introduction}  %Title of the First Chapter

\ifpdf
    \graphicspath{{Chapter1/Figs/Raster/}{Chapter1/Figs/PDF/}{Chapter1/Figs/}}
\else
    \graphicspath{{Chapter1/Figs/Vector/}{Chapter1/Figs/}}
\fi


%********************************** %First Section  **************************************
Elections are ubiquitous element of our society in one form or another.
For example, choosing members of parliament or a winner of a sports competition,
in day-to-day activities, we need to schedule a meeting time, or a video streaming platform needs to provide recommended movies list to the user.
Each of these choices may require different type of election and different strategy of selecting the winner (or winners).
The differences come from the underlying properties and objectives of the choice.
For example, in the presidential election, we want to choose the candidate supported by the majority of the voters.
It is not the case for the parliamentary elections, where the goal is to select members of the parliament in such a way that 
the opinions of the voters are represented proportionally. 
\textcolor{orange}{(WIP: Czy potrzeba jeszcze wiecej przykładów, np. z , facility location, blockchain, economics...)}

More formally, a committee election is the process of choosing a fixed-size subset of candidates (committee), 
based on the voters' preferences. Voters' ballots (preferences) can be represented in various different ways.
In this work we are going to focus on approval elections, that is,
on elections where each voter provides a subset of candidates that he or she approves of.
Another notable type of committee elections is ordinal model,
where the voters order the candidates from the most to the least preferred. 
See the work of \cite{faliszewskiSkowronMultiWinnerOverview},
which provides the overview and discusses various committee election types,
and for the extensive survey of approval model and current state of the research in this area, we point to \cite{lacknerSkowronAVOverview}. 

Winner determination is one of the core problems in computational social choice theory, however not the only one.
We need to mention control problems, it refers to certain ways of manipulating elections.
In short, we try to make a given candidate win (or lose), by changing the structure of an election,
that is adding or deleting voters or candidates.
Other interesting problems, that we are not going to explore in this work,
are manipulation, bribery, justified representation and many others. \textcolor{orange}{(WIP: Cytacje?)}

Many of the mentioned above social choice problems are known to be computationally hard.
We can deal with this complexity using various methods,
more standard ones, like approximate (random or deterministic) or parametrized algorithms.
An alternative approach to these classical solutions is introducing restrictions on preference domains,
in other words, we restrict the structure of voters ballots (preferences). 
There is a successful and active line of research in this area.
Beginning with a notion of \textit{single-peaked} introduced by \cite{Black-single-peaked} 
and \textit{single-crossing} by \cite{mirrlees-single-crossing}.
Arriving to much more recent work, \cite{DBLP:journals/corr/ElkindL15} explored well known restrictions in the setting of approval preferences, and
introduced several new one-dimensional domain restriction models such as \textit{Voter Interval (VI)} and \textit{Candidate Interval (CI)},
both based on the \textit{Dichotomous Euclidean} restriction. 
\cite{falisz} improved this model by making it more versatile,
due to using additional radius parameter and providing N-dimensional definition.
\textcolor{orange}{(WIP: Jaki był aspekt uzupełnienia modeli euklidesowych? Czy tu jest za wcześnie na mówienie o promieniu?)}
Calling the new model \textit{Voter Candidate Range (VCR)}, and later they analyzed two-dimensional case and briefly one-dimensional,
for committee election problems.

\subsubsection{\textcolor{orange}{(WIP: Bez tytułu?)} Contributions}
This work's main contributions are following:
\begin{enumerate}
    \item We provide detection algorithm for VCR property, based on the ILP. 
    Additionally, we extended it to detect VR and CR properties, and compare it with existing solutions. 
    \item We analyze VCR profiles domain extensively, and compare it to other well known models.
    \item We propose control problem polynomial time algorithm for the VCR model,
    expanding on the solutions for CR and VR models,
    from \cite{control-cr-how-hard-magiera-faliszewski} and \cite{control-sp-shield-faliszewski-hemaspaandra} respectively.
    \item \textcolor{orange}{(WIP: Napisać trochę więcej i podać numery rozdziałów/sekcji)}
\end{enumerate}

We do not discuss the related work right now, instead we do it throughout this thesis,
as the discussion will greatly benefit from the added context and introduced definitions in the subsequent chapters.